% Folder 91, batch 07 — archivists' pages 121–140.
% Pass: Opus 5.5 (claude-opus-5-5), 2026-10-03 — first pass, unchecked against the pages by a human.
%
% Pages 122, 126, 128, 138 and 140 are typed versos unrelated to the notes;
% pages 124, 129, 132 and 136 are blank. They are absent.
\documentclass[11pt,a4paper]{article}
\input{../preamble/grothendieck}

\folder{91}
\batch{7}
\pages{121}{140}
\foldertitle{Autour de Néron / Greenberg-Néron. Foncteurs Hom (méthodes non-projectives) : notes manuscrites (s.d.), lettre (1967).}
\dating{[à partir de 1964-vers 1970]}
\watermark{Édition de démonstration}

\begin{document}

\note{la p. 121 reprend un énoncé commencé avant le lot : les hypothèses sur $X$ et la numérotation (i)–(iii) s'y enchaînent sans titre.}

\page{121}
$X$ \uncertain{schéma} \uncertain{compact}, [i.e. \uncertain{quasi-compact} et \uncertain{séparé}\add{\uncertain{topologiquement}}]

(i) $X$ est tot. discontinu, \struck{i.e. si $x, y \in X$, $x \neq y$, alors il existe une partie ouverte et fermée de $X$ qui contient $x$ et non $y$. En effet, \ill{} $U$, $V$ des ouverts, $U \ni x$, $V \ni y$, $U \cap V = \emptyset$.} i.e. tout $x \in X$ a un syst. fond. de voisinages \uncertain{qui sont} à la fois ouverts et fermés.

[En fait, il y a un syst. fond. de voisinages ouverts quasi-compacts]

(ii) $X$ \struck{est} affine. En effet, \add{on recouvre $X$ par des} \struck{\ill{} \ill{}} ouverts affines \struck{\ill{}} finis, soit $X_i$ ($1 \leq i \leq n$) \struck{affines \ill{} \ill{} qui se \ill{}} \struck{\ill{}}. Ils sont ouverts et fermés. Posons
\[
  X'_i = X_i - X_i \cap \bigcup_{j<i} X_j ,
\]
Alors $X'_i$ est une partie ouverte et fermée de $X_i$, donc affine \uncertain{puisque} $X_i$ l'est (un ouvert \uncertain{p. ex.} \ill{} \ill{} \ill{} \ill{} \ill{}). Donc $X$ \struck{\ill{}} est somme des $X'_i$, il est affine.

(iii) Les parties constructibles de $X$ sont les \uncertain{parties} ouvertes et fermées. [\struck{Donc} \add{\uncertain{En effet}} car $X \to x$ est un \struck{\ill{}} \uncertain{bon}] Donc
\note{la page s'arrête sur « Donc » ; la suite n'est pas dans le lot.}

\page{123}
\emph{Greenberg -- Néron}
\note{seul mot de la feuille, en haut à droite : un titre de sa main pour les feuillets qui suivent.}

\page{125}
Soit $\mathfrak{X} = (X_i)_{i \in I}$ un système projectif de préschémas, ensemblistement \uncertain{stationnaire} [i.e. $\exists\, i_0$, $\forall\, i \geq i_0$, \ill{} \ill{}]. Alors $X = \varprojlim_i X_i$ existe, $\operatorname{Top} X = \varprojlim_i \operatorname{Top}(X_i)$ ---.

Soit $x \in X$, on dit que $x$ est \emph{épais} si

(Ep) $\exists\, i$ t.q. $\forall\, j \geq i$, $x_j$ est \uncertain{maximal} dans $\varphi_{ij}^{-1}(x_i)$, donc si $k \geq j \geq i$, $x_k$ est \uncertain{maximal} dans $\varphi_{jk}^{-1}(x_j)$.

Sous quelles conditions est-ce équivalent à :

(Ep)$'$ $\exists\, i$, \ill{} un voisinage ouvert non vide $U_i$ de $\bar{x}_i$, t.q. $\forall\, j \geq i$, $\varphi_{ij}^{-1}(U_i) \subset \bar{x}_j$ (donc la même chose pour \uncertain{tout} $i' \geq i$).

$=$ propriété \ill{} \ill{} \ill{} : $\varphi_{ij}^{-1}(U_i)$ est irréductible, de pt générique $x_j$, et donc (Ep)$'$ implique (Ep) !

Conditions suffisantes :

(Ep)$''$ $\bar{x}$ \marginal{entouré : $\bar{x} = \varprojlim \bar{x}_i$} est une partie constructible de $X$. \add{(si les $X_i$ sont noethériens)}

Ceci équivaut à : $\bar{x} = \varphi_i^{-1}(\bar{x}_i)$ pour $i$ grand

\struck{Cons. : pour $i$ \ill{} constructible \ill{}} (dès que $i' \geq i$)

(Ep)$'''$ $\exists$ une partie constructible \add{ouverte non vide} \struck{\ill{}} de $\bar{x}$, \uncertain{soit} $U$. i.e. $\exists\, U_i$
\note{la phrase s'arrête là ; elle reprend en haut de la p. 127 (la p. 126 est étrangère).}

\page{127}
$x_i$ \uncertain{maximal} dans $\varphi_{ji}^{-1}($\ill{}$)$

\note{encadré de sa main ; les implications sont des flèches doubles. Transcrit en liste.}
\begin{itemize}
  \item (Ep)$'$ $\Longrightarrow$ (Ep) $+$ (\uncertain{unicité} : \ill{} \ill{} de $X$) ;
  \item (Ep)$''$ $\Longrightarrow$ (Ep)$'$, la flèche portant « M.L. et $\Sigma = \mathbf{N}$ » ;
  \item (Ep)$''$ $\Longrightarrow$ (Ep)$'''$, la flèche portant « 4 » ; \struck{(Ep)$'''$} ;
  \item (Ep)$'''$ $\Longrightarrow$ (Ep)$'$ si \struck{$\forall x$} \ill{} $X'_i$, $\exists$ un ouvert $\neq \emptyset$ $U_i \subset \bar{x}_i$ t.q. les $\varphi_{ij}^{-1}(U_i) \times X'_j = \varphi_j(X)$ \quad ???
\end{itemize}

\ill{} pt épais $\longleftarrow$ parties constructibles \add{irréductibles}

\[
  U_j \cap X'_j \subset \bar{x}_j \qquad\qquad X'_i \longleftarrow X'_j
\]
\[
  U_i \subset \bar{x}_i
\]
\[
  x \in Z \subset \bar{x} \qquad x \in U_0 \cap Z_0 \qquad \bar{x}
\]
$x \in$ \struck{$\bar{x}$} $\bigcup_i U_i \cap Z_i$

$x \in U_0 \cap Z_0 \subset \bar{x}$ \qquad $\bar{x} = \overline{U_0 \cap Z_0} \subset Z_0$

\struck{$x \in \ldots \in Z_0$}

\uncertain{Donc} $f_1 \ldots f_n$ \ill{} \ill{} $x$

$f_1 \ldots f_n \in \mathfrak{p}$ \qquad $g$

si $g \notin \mathfrak{q}$ \qquad $\mathfrak{q} \supset \mathfrak{p} \Longleftrightarrow \mathfrak{q} \supset f_1 \ldots f_n$

\drawing{deux flèches descendent de $\mathfrak{X}$ et de $\mathfrak{X}_\xi$ et convergent vers un même point, sans but écrit.}

$\mathcal{O}_{X,x}$ \qquad $\mathfrak{m}_x$ \qquad \ill{}
\note{au bas de la page, autour de $\mathcal{O}_{X,x}$ et $\mathfrak{m}_x$, quelques signes épars, illisibles.}

\page{130}
$V$ anneau de valuation discrète, corps résiduel $k$, parfait, corps des fractions $K$.

Pour tout préschéma $X$ sur $V$, on définit $\bar{X}$ \uncertain{par} Greenberg.
\[
  \bar{X}(k') \simeq X(V_{k'}) \qquad \text{\struck{$\to X$}}
\]
($k'$ alg. sur $k$, $V_{k'}$ alg. \uncertain{correspondante} \uncertain{étale} sur $V$)

Si $k'$ est parfait, $V_{k'}$ est un anneau de valuation discrète, \uncertain{si} $K_{k'} = K'$ son corps des fractions, \ill{}.
\[
  \bar{X}(k') \simeq X(V_{k'}) \longrightarrow X(K')
\]
Soit \struck{\ill{}} $\to Y$ un \uncertain{morphisme} de $V$-préschémas, \add{\uncertain{propre} sur $V$,} et soit
\[
  \varphi : X_K \longrightarrow Y_K
\]
un morphisme de $K$-préschémas, \ill{} \ill{} un \uncertain{morphisme} \ill{} \ill{}

\marginal{en biais : si $Y$ n'est pas propre sur $V$, \ill{} \ill{} \ill{} \ill{} « application \ill{} » \ill{} $\bar{X} \to \bar{Y}$, ayant des propriétés \ill{}}

On définit une application \emph{ensembliste}
\[
  \bar{\varphi} : \bar{X} \longrightarrow \bar{Y}
\]
Soit en effet $\bar{x} \in \bar{X}$, soit $k' = \tilde{k}(\bar{x})$ « la » clôture parfaite de $k(\bar{x})$, alors $\bar{x}$ définit un pt $\bar{x}' \in \bar{X}(k')$, d'où un pt, \struck{\ill{}} $\alpha_X(\bar{x}') \in X(V_{k'}) \subset X(K_{k'})$ et \uncertain{appliquant} $\varphi$ un point
\[
  \varphi(\alpha_X(\bar{x}')) \in Y(K_{k'}) \xleftarrow{\ \sim\ } Y(V_{k'}) \overset{\alpha_Y}{\simeq} \bar{Y}(k')
\]
\note{le $Y$ de $Y(K_{k'})$ porte des traits au-dessus (barre et flèches) ; lecture incertaine.}
donc la \uncertain{localité} \ill{} \ill{} $\bar{\varphi}(x)$. On a \ill{} \ill{} par construction un \ill{} \uncertain{homomorphisme}
\[
  k(\varphi(\bar{x})) \longrightarrow \widetilde{k(\bar{x})}
\]
(\uncertain{où} $\widetilde{\phantom{k}}$ désigne la clôture \add{\uncertain{séparable}} parfaite).

Par la suite, $X$, $Y$ de type fini sur $K$, et \add{\uncertain{séparés}} intègres.

\emph{Proposition 1} Supposons $\bar{x} \in \bar{X}$ \emph{épais} et $\varphi$ \struck{\emph{dominant} \uncertain{génér.}t} \add{\uncertain{excité}, et} \add{\uncertain{en $x$}} \emph{séparable} [i.e. \ill{} \ill{} \ill{}]. Alors \add{(\uncertain{si})} $\bar{y} = \varphi(\bar{x})$ est \emph{épais}, $k(\bar{y}) \subset k(\bar{x})$, et l'application \uncertain{ponctuelle} $\bar{x}_{\mathrm{red}} \to \bar{y}_{\mathrm{red}}$ définie \ill{} provient d'un morphisme d'indice fini \uncertain{constructible} dans un $U \in \bar{x}$, [d'image un \ill{} \add{ouvert} constructible de $\bar{y}_{\mathrm{red}}$].
\note{l'énoncé est marqué d'un trait vertical en marge gauche ; la lecture de « dominant génér.t » et de ses corrections interlinéaires est incertaine.}

\page{131}
\emph{Corollaire} Fonctorialité \add{(en $X_K$)} de l'ensemble $\operatorname{Ép}(\bar{X})$ des points \struck{\ill{}} épais de $\bar{X}$ [\ill{} \ill{} \ill{} \ill{} corps résiduels ----], par des \add{\emph{applications rationnelles}} \struck{\emph{morphismes}} $X_K \to Y_K$ \emph{dominantes génér.t séparables} de $K$-\emph{schémas propres}

\emph{Proposition 2} Soit $S$ une partie constructible de $\bar{X}$, $X_K$ \uncertain{intègre}, [tel que $(V, x)$ y soit bonne], et tel que $\bar{\varphi}(S)$ existe [\uncertain{ex.} $Y$ propre sur $V$] et soit $Y_K$-\uncertain{simple}. Alors $\bar{\varphi}(S)$ est constructible. De plus, il y a une \add{\uncertain{bonne}} partition $S = \coprod_i S_i$, les $S_i$ des sous-schémas irréductibles constructibles de $X$, telles que $\varphi$ soit \uncertain{morphique} d'indice fini sur chaque $S_i$, [et $\varphi(S_i)$ constructible irréductible \ldots]
\ill{} alors \ill{} \ill{} \ill{}

\emph{Corollaire} Fonctorialité \struck{de} en $X_K$ variant comme dessus, et de plus simple sur $K$, de l'ens $\operatorname{Cons}(\bar{X})$ des parties constructibles de $\bar{X}$.

\page{133}
Condition de \uncertain{finitude} d'anneaux [ou de préschémas \uncertain{quasi-compacts}]
\[
  \Gamma : A \rightsquigarrow A((t))^{*} = \bigl(A[[t]]_t\bigr)^{*} \qquad (t \text{ une indéterminée})
\]
Ce foncteur est ind-représentable par un
\[
  \mathcal{G} = \varinjlim_n (\mathcal{G}_n)
\]
\marginal{$\varinjlim$ par \ill{} de \ill{}}
où $\mathcal{G}_n$ représente le sous-foncteur $\Gamma_n$ de $\Gamma$, \ill{}
\[
  \text{\struck{$A \rightsquigarrow A((t))^{*}$}}
\]
\uncertain{formé} des séries formelles $\sum_{-\infty}^{+\infty} a_i t^i$ \add{\uncertain{qui ont un inverse} $\sum_{-\infty}^{+\infty} b_j t^j$, \ill{} dans $\Gamma_n$ \ill{}} $a_i = 0$, $b_i = 0$ pour $i < -n$. $\mathcal{G}_n$ est \uncertain{affine} \ill{} \ill{} un sous-schéma fermé de $\operatorname{Spec} \mathbf{Z}[(a_i, b_i)_{i \geq -n}]$, [\uncertain{défini} par les relations \uncertain{exprimant}
\[
  \Bigl(\sum_i a_i t^i\Bigr)\Bigl(\sum_j b_j t^j\Bigr) = 1 ,
\]
savoir $a_{-n} b_{-n} = 0$, … ]

\struck{En \ill{}} Soit, $\Gamma_0$ est le foncteur
\[
  \Gamma_0 : A \rightsquigarrow A[[t]]^{*}
\]
qui est représenté par
\[
  \mathcal{G}_0 = \mathbf{G}_m \times \operatorname{Spec} \mathbf{Z}[(a_i)_{i \geq 1}]
\]
Considérons le schéma
\[
  \mathcal{G}' = \mathbf{Z}_{\mathbf{Z}} \times \mathcal{G}_0
\]
qui représente un foncteur ind-représentable
\[
  \mathcal{G}' = \varinjlim_n \mathcal{G}'_n , \qquad \mathcal{G}'_n = [-n, n]_{\mathbf{Z}} \times \mathcal{G}_0 .
\]
On a un morphisme canonique
\[
  \mathcal{G}' \longrightarrow \mathcal{G}
\]

\page{134}
correspondant : un morphisme
\[
  \mathcal{G}'_n \longrightarrow \mathcal{G}_n
\]
\struck{\ill{} \ill{} \ill{} \uncertain{famille}} \ill{} \ill{}, \uncertain{correspondant} \uncertain{par} \uncertain{composantes}, à
\[
  \varphi_n^{\nu} : \mathcal{G}_0 \longrightarrow \mathcal{G}_n \qquad \nu \in [-n, +n]
\]
donnés \struck{par} sur les \uncertain{foncteurs} par
\[
  \varphi_n^{\nu}\Bigl(\text{\struck{$\Sigma$}}\, a_i t^i_{\ i \geq 0}\Bigr) = t^{\nu} \sum a_i t^i
\]
Ceci posé, on a

\emph{Th} (i) Le morph. $\varphi_n : \mathcal{G}'_n \to \mathcal{G}_n$ est une \uncertain{immersion} \emph{fermée} \emph{surjective}.

(ii) Elle induit une immersion ouverte sur $\mathcal{G}_0 \times \{\nu\}$ pour $\nu = \pm n$, mais pour $-n < \nu < +n$, le morphisme induit n'est \emph{pas} une \uncertain{immersion} [\ill{} \uncertain{même} \uncertain{ouverte} \ill{} \ill{} \ill{} \uncertain{dès} \ill{} $S \to \operatorname{Spec} \mathbf{Z}$, $S \neq \emptyset$].

(iii) Soit $\mathcal{G}_n^i$ le sous-schéma ouvert \add{\uncertain{induit}} et fermé de $\mathcal{G}_n$ \struck{\ill{}} sur l'image ensembliste \struck{\ill{}} $\varphi_n(\mathcal{G}_0 \times \{i\})$, avec $|i| \leq n$.
\marginal{en biais : donc $\mathcal{G}_n$ est réunion disjointe d'ouverts $\mathcal{G}_n^i$ ($i \in [-n, +n]$), d'ailleurs isomorphes aux images des $\varphi_n^i$}

\uncertain{Pour} $n$ \ill{}, \ill{} \ill{} \ill{} un système \struck{\ill{}} \add{\uncertain{d'immersions}} \uncertain{inductif} de préschémas. \ill{} $n' \geq n \geq |i|$, les morphismes
\[
  \mathcal{G}_n^i \longrightarrow \mathcal{G}_{n'}^i
\]
\uncertain{induits} \ill{} \ill{} \uncertain{isomorphismes}, et \ill{} \ill{} \ill{} \ill{} $S \neq \emptyset$ \uncertain{quelc.} [\ill{} \ill{} \ill{} \ill{} \ill{} corps \uncertain{premiers} …].

\page{135}
\note{le schéma ci-dessous redresse le dessin : les signes $\cup$ entre les lignes (« sous-catégorie pleine ») sont rendus par des inclusions, et chaque catégorie est glosée à droite.}

\begin{tikzcd}[column sep=small, row sep=large, nodes={font=\scriptsize}]
 & \operatorname{Pro}(\mathrm{Sch}_{/S}) \\
 (\mathrm{Sch}_{/S}) & \operatorname{Pro}'(\mathrm{Sch}_{/S}) \arrow[l, "\varprojlim"'] \arrow[u, hook, "\text{sous-catégorie pleine}"'] \\
 & \operatorname{Pro}''(\mathrm{Sch}_{/S}) \arrow[u, hook, "\text{sous-catégorie pleine}"'] \arrow[ul, "\text{foncteur pleinement fidèle}"]
\end{tikzcd}

$\operatorname{Pro}(\mathrm{Sch}_{/S})$ : pro-préschémas

$\operatorname{Pro}'(\mathrm{Sch}_{/S})$ : pro-préschémas $(X_i)_{i \in I}$, t.q. $\exists\, i$, $u_{ji}$ affine pour $j \geq i$

$\operatorname{Pro}''(\mathrm{Sch}_{/S})$ : pro-préschémas $(X_i)_{i \in I}$ t.q.
\[
  \begin{cases}
    \forall\, i \text{ grand}, X_i \text{ est qu.-compact et de présentation finie sur } S \\
    \exists\, i \text{ t.q. } j \geq i \text{ implique } u_{ji} \text{ affine}
  \end{cases}
\]
\note{au-dessus, une première rédaction biffée : \struck{$\exists\, i$, [$X_i$ \ill{} qu.-compact et \ill{} présentation finie ; $u_{ji}$ affine de prés. finie]}.}

Soient $\mathfrak{X} \in \operatorname{Ob} \operatorname{Pro}''(\mathrm{Sch}_{/S})$, $\mathcal{Y} \in \operatorname{Ob} \operatorname{Pro}'(\mathrm{Sch}_{/S})$ \uncertain{alors} \add{($\mathcal{Y} = (Y_j)$, les $Y_j$ quasi-compacts)}
\[
  \operatorname{Hom}_{\operatorname{Pro}(\mathrm{Sch}_{/S})}(\mathcal{Y}, \mathfrak{X}) \xrightarrow{\ \sim\ } \operatorname{Hom}\bigl(\varprojlim \mathcal{Y}, \varprojlim \mathfrak{X}\bigr) ,
\]
\struck{[En particulier, la catégorie} \add{\uncertain{Si} \ill{} \ill{} \ill{} base $S$}
\[
  \operatorname{Hom}_{\operatorname{Pro}(\mathrm{Sch}_{/S})}(\mathcal{Y}, \mathfrak{X}) \simeq \operatorname{Hom}_S\bigl(\varprojlim \mathcal{Y}, \varprojlim \mathfrak{X}\bigr)
\]
\struck{Dans les propriétés de Serre}

\emph{Remarque} Dans Serre et Néron, on utilise des \add{(\uncertain{compléments})} systèmes projectifs de schémas \add{(de \uncertain{particuliers} \uncertain{types})} qui sont $\in \operatorname{Pro}''(\mathrm{Sch})$. \struck{D} Il y a avantage à les « interpréter » comme étant de véritables schémas sur $k$, \ill{} au lieu de pro-objets ……

\page{137}
\struck{$a_{i0}$} \qquad $a_i, b_{ij}, c_{ip} \in A$
\[
  \begin{cases}
    f_i = a_i + \displaystyle\sum_{1 \leq j \leq n} b_{ij}\, t_j + \sum_{|p| \geq 2} c_{ip}\, t^p \\
    1 \leq i \leq n
  \end{cases}
  \qquad
  \begin{cases}
    a_i \in \mathfrak{m} \\
    \det(b_{ij}) \notin \mathfrak{m}
  \end{cases}
\]
\[
  (\lambda_1, \ldots, \lambda_n) \in \mathfrak{m}^{(n)} \xmapsto{\ \lambda_{(f = (f_1, \ldots, f_n))}\ } \bigl(f_1(\lambda_1, \ldots, \lambda_n), \ldots, f_n(\lambda_1, \ldots, \lambda_n)\bigr)
\]
Si $A$ est d'\uncertain{échelon} \uncertain{fini} $N$, \uncertain{alors} $\lambda_f$ est \uncertain{connu} quand on connaît $f \bmod (t_1, \ldots, t_n)^{N+1}$, donc quand on connaît \struck{les} \struck{\uncertain{polynômes}} \uncertain{les} \ill{} \uncertain{de} \uncertain{degré} $\leq N$.
\[
  \boxed{\, f_i(\lambda_1, \ldots, \lambda_n) = a_i + \sum_j b_{ij}\, \lambda_j + \sum_{2 \leq |p| \leq N} c_{ip}\, \lambda_1^{p_1} \cdots \lambda_n^{p_n} \,}
\]
\[
  \begin{array}{c}
    B \simeq A[[t_1, \ldots, t_n]] \\
    \uparrow \\
    A
  \end{array}
  \qquad
  \begin{array}{l}
    V_1 \; [\; \mathfrak{m}^{?},\ \mathfrak{m}^2 \\
    V_2 \; [\; \mathfrak{m}^3 \\
    \quad\ \vdots \\
    \phantom{V_N} \; [\; \mathfrak{m}^{N_1},\ \mathfrak{m}^{N_?} \\
    V_N \; [\; \mathfrak{m}^{N_?} = 0
  \end{array}
\]
\note{à droite du diagramme $B \leftarrow A$, une échelle de puissances de $\mathfrak{m}$ découpée en tranches $V_1, \ldots, V_N$ par des crochets ; les exposants notés « ? » sont illisibles, seule la fin $= 0$ est sûre.}

\page{139}
\[
  G_{n,N} = \underline{\operatorname{Aut}}_{\text{algèbres augmentées}} \mathbf{Z}[t_1, \ldots, t_n]/(t_1, \ldots, t_n)^{N+1}
\]
$\Gamma_n = \operatorname{Grb}(G_{n,N})$ groupe algébrique sur $k$ des changements de coordonnées
\[
  \begin{cases}
    f_i(t_1, \ldots, t_n) = \displaystyle\sum_j b_{ij}\, t_j + \sum_{2 \leq |p| \leq N} c_{ip}\, t_1^{p_1} \cdots t_n^{p_n} \\
    b_{ij}, c_{ip} \in A \qquad \det(b_{ij}) \notin \mathfrak{m}
  \end{cases}
\]
\note{l'indice de $\Gamma_n$ est surchargé d'une tache ; « $\operatorname{Grb}$ » est lu comme le foncteur de Greenberg.}

Ce groupe opère sur $\mathfrak{m}^n \simeq \underline{A}^n$ \uncertain{si} \uncertain{l'échelon} \uncertain{infinitésimal} \uncertain{est} $\leq N$. \uncertain{Donc}

$g \in \Gamma_n$, $a \in \mathfrak{m}^n$ \qquad $x \mapsto T_g(x) \dotplus a$
\[
  \text{\struck{$T_g$}}\ b \dotplus T_g\bigl(\text{\struck{$T$}}\, a \dotplus T_f(x)\bigr)
\]
\[
  b \dotplus T_g(a) \dotplus
\]
\[
  T_{g'}(a + b) \qquad\qquad g(x + y) = g(x)
\]
\[
  \text{\struck{$T_i$}} \qquad \text{\struck{$g$}}\ g(a + b) = g(a) + g'
\]
\[
  g(x + y) =
\]
\note{le calcul de la moitié inférieure s'interrompt sans conclusion ; « $\dotplus$ » rend une addition marquée d'un point au-dessus, lecture incertaine.}

\end{document}
